\documentclass[12pt,a4paper]{article}

\usepackage{babel}
\usepackage[utf8]{inputenc}
\usepackage[T1]{fontenc}
\usepackage[margin=2.2cm]{geometry}
\usepackage{hyperref}
\usepackage{enumitem}
\usepackage{booktabs}
\usepackage{xcolor}
\usepackage{titlesec}
\usepackage{titling}
\usepackage{microtype}
\usepackage{array}
\usepackage{multirow}
\usepackage{listings}
\usepackage{tcolorbox}
\usepackage{float}
\usepackage{lmodern} % fuente segura para pdfLaTeX

% Colores
\definecolor{primary}{RGB}{15,118,110}
\definecolor{secondary}{RGB}{37,99,235}
\definecolor{muted}{RGB}{107,114,128}
\definecolor{bgcode}{RGB}{248,250,252}
\definecolor{frame}{RGB}{226,232,240}

% Estilo de secciones
\titleformat{\section}{\Large\bfseries\color{primary}}{\thesection}{0.6em}{}
\titleformat{\subsection}{\large\bfseries\color{secondary}}{\thesubsection}{0.6em}{}
\titleformat{\subsubsection}{\bfseries\color{secondary}}{\thesubsubsection}{0.6em}{}

% Listings (SQL)
\lstdefinelanguage{SQLPlus}[]{SQL}{
  morekeywords={IF,ELSE,END,DECLARE,BEGIN,EXCEPTION,WHEN,LOOP,ELSIF,THEN,RAISE,RETURN},
  sensitive=false
}
\lstset{
  language=SQLPlus,
  basicstyle=\ttfamily\small,
  keywordstyle=\color{secondary}\bfseries,
  commentstyle=\color{muted}\itshape,
  stringstyle=\color{primary},
  columns=fullflexible,
  backgroundcolor=\color{bgcode},
  frame=single,
  rulecolor=\color{frame},
  breaklines=true,
  tabsize=2,
  showstringspaces=false,
    inputencoding = utf8,  % Input encoding
    extendedchars = true,  % Extended ASCII
    literate      =        % Support additional characters
      {á}{{\'a}}1  {é}{{\'e}}1  {í}{{\'i}}1 {ó}{{\'o}}1  {ú}{{\'u}}1
      {Á}{{\'A}}1  {É}{{\'E}}1  {Í}{{\'I}}1 {Ó}{{\'O}}1  {Ú}{{\'U}}1
      {à}{{\`a}}1  {è}{{\`e}}1  {ì}{{\`i}}1 {ò}{{\`o}}1  {ù}{{\`u}}1
      {À}{{\`A}}1  {È}{{\`E}}1  {Ì}{{\`I}}1 {Ò}{{\`O}}1  {Ù}{{\`U}}1
      {ä}{{\"a}}1  {ë}{{\"e}}1  {ï}{{\"i}}1 {ö}{{\"o}}1  {ü}{{\"u}}1
      {Ä}{{\"A}}1  {Ë}{{\"E}}1  {Ï}{{\"I}}1 {Ö}{{\"O}}1  {Ü}{{\"U}}1
      {â}{{\^a}}1  {ê}{{\^e}}1  {î}{{\^i}}1 {ô}{{\^o}}1  {û}{{\^u}}1
      {Â}{{\^A}}1  {Ê}{{\^E}}1  {Î}{{\^I}}1 {Ô}{{\^O}}1  {Û}{{\^U}}1
      {œ}{{\oe}}1  {Œ}{{\OE}}1  {æ}{{\ae}}1 {Æ}{{\AE}}1  {ß}{{\ss}}1
      {ẞ}{{\SS}}1  {ç}{{\c{c}}}1 {Ç}{{\c{C}}}1 {ø}{{\o}}1  {Ø}{{\O}}1
      {å}{{\aa}}1  {Å}{{\AA}}1  {ã}{{\~a}}1  {õ}{{\~o}}1 {Ã}{{\~A}}1
      {Õ}{{\~O}}1  {ñ}{{\~n}}1  {Ñ}{{\~N}}1  {¿}{{?`}}1  {¡}{{!`}}1
      {„}{\quotedblbase}1 {“}{\textquotedblleft}1 {–}{$-$}1
      {°}{{\textdegree}}1 {º}{{\textordmasculine}}1 {ª}{{\textordfeminine}}1
      {£}{{\pounds}}1  {©}{{\copyright}}1  {®}{{\textregistered}}1
      {«}{{\guillemotleft}}1  {»}{{\guillemotright}}1  {Ð}{{\DH}}1  {ð}{{\dh}}1
      {Ý}{{\'Y}}1    {ý}{{\'y}}1    {Þ}{{\TH}}1    {þ}{{\th}}1    {Ă}{{\u{A}}}1
      {ă}{{\u{a}}}1  {Ą}{{\k{A}}}1  {ą}{{\k{a}}}1  {Ć}{{\'C}}1    {ć}{{\'c}}1
      {Č}{{\v{C}}}1  {č}{{\v{c}}}1  {Ď}{{\v{D}}}1  {ď}{{\v{d}}}1  {Đ}{{\DJ}}1
      {đ}{{\dj}}1    {Ė}{{\.{E}}}1  {ė}{{\.{e}}}1  {Ę}{{\k{E}}}1  {ę}{{\k{e}}}1
      {Ě}{{\v{E}}}1  {ě}{{\v{e}}}1  {Ğ}{{\u{G}}}1  {ğ}{{\u{g}}}1  {Ĩ}{{\~I}}1
      {ĩ}{{\~\i}}1   {Į}{{\k{I}}}1  {į}{{\k{i}}}1  {İ}{{\.{I}}}1  {ı}{{\i}}1
      {Ĺ}{{\'L}}1    {ĺ}{{\'l}}1    {Ľ}{{\v{L}}}1  {ľ}{{\v{l}}}1  {Ł}{{\L{}}}1
      {ł}{{\l{}}}1   {Ń}{{\'N}}1    {ń}{{\'n}}1    {Ň}{{\v{N}}}1  {ň}{{\v{n}}}1
      {Ő}{{\H{O}}}1  {ő}{{\H{o}}}1  {Ŕ}{{\'{R}}}1  {ŕ}{{\'{r}}}1  {Ř}{{\v{R}}}1
      {ř}{{\v{r}}}1  {Ś}{{\'S}}1    {ś}{{\'s}}1    {Ş}{{\c{S}}}1  {ş}{{\c{s}}}1
      {Š}{{\v{S}}}1  {š}{{\v{s}}}1  {Ť}{{\v{T}}}1  {ť}{{\v{t}}}1  {Ũ}{{\~U}}1
      {ũ}{{\~u}}1    {Ū}{{\={U}}}1  {ū}{{\={u}}}1  {Ů}{{\r{U}}}1  {ů}{{\r{u}}}1
      {Ű}{{\H{U}}}1  {ű}{{\H{u}}}1  {Ų}{{\k{U}}}1  {ų}{{\k{u}}}1  {Ź}{{\'Z}}1
      {ź}{{\'z}}1    {Ż}{{\.Z}}1    {ż}{{\.z}}1    {Ž}{{\v{Z}}}1  {ž}{{\v{z}}}1
}

% Caja de aviso
\tcbset{
  colback=white,
  colframe=primary,
  coltitle=white,
  colbacktitle=primary,
  fonttitle=\bfseries,
  arc=2mm,
  boxrule=0.6pt
}


\title{\textbf{Guía de Auditoría y Control de Datos en MySQL}}
\author{}
\date{}

\begin{document}

\maketitle

Para garantizar que el despliegue de permisos es exacto y evaluar de forma eficiente el trabajo en la terminal, no basta con ejecutar las sentencias de control de datos (DCL). Es fundamental auditar el diccionario de datos interno de MySQL.

A continuación, se presenta la guía completa estructurada, incorporando las consultas de verificación directas a las tablas del sistema para comprobar que cada paso se ha ejecutado correctamente.

\section{Arquitectura de Cuentas: Usuario y Host}
En MySQL, una cuenta se define estrictamente por la combinación de nombre de usuario y el host desde el que se conecta (\texttt{'usuario'@'host'}).

\subsection{Creación y Modificación}
\textbf{Crear un usuario:}
\begin{lstlisting}
CREATE USER 'dev_app'@'10.0.0.%' IDENTIFIED BY 'PasswordSegura123!';
\end{lstlisting}

\textbf{Bloquear una cuenta temporalmente:}
\begin{lstlisting}
ALTER USER 'dev_app'@'10.0.0.%' ACCOUNT LOCK;
\end{lstlisting}

\subsection*{ Verificación de Cuentas}
Para comprobar que el usuario se ha creado con los atributos correctos, consulta la tabla principal del sistema:
\begin{lstlisting}
SELECT host, user, plugin, authentication_string, account_locked 
FROM mysql.user 
WHERE user = 'dev_app';
\end{lstlisting}

\section{Gestión de Privilegios (GRANT / REVOKE)}
La asignación debe basarse en el principio de mínimo privilegio. Evita los permisos globales (\texttt{*.*}) y ajusta el acceso a la base de datos o tabla específica.

\subsection{Asignación y Revocación}
\textbf{Nivel de Base de Datos:}
\begin{lstlisting}
GRANT SELECT, INSERT, UPDATE, DELETE ON tienda_web.* TO 'app_backend'@'localhost';
\end{lstlisting}

\textbf{Nivel de Tabla y Columna:}
\begin{lstlisting}
GRANT SELECT, UPDATE (precio, stock) ON tienda_web.productos TO 'gestor_precios'@'%';
\end{lstlisting}

\textbf{Revocar permisos:}
\begin{lstlisting}
REVOKE DELETE ON tienda_web.* FROM 'app_backend'@'localhost';
\end{lstlisting}

\subsection*{ Verificación de Privilegios}
Aunque \texttt{SHOW GRANTS FOR 'usuario'@'host';} es útil, para automatizar comprobaciones es mejor consultar las tablas de privilegios específicas:

\begin{lstlisting}
-- Verificar privilegios globales (deben estar en 'N' para usuarios normales)
SELECT user, host, Select_priv, Insert_priv, Super_priv
FROM mysql.user WHERE user = 'app_backend';

-- Verificar privilegios a nivel de base de datos
SELECT host, user, db, Select_priv, Insert_priv
FROM mysql.db WHERE user = 'app_backend';

-- Verificar privilegios a nivel de tabla
SELECT host, user, db, table_name, table_priv, column_priv
FROM mysql.tables_priv WHERE user = 'gestor_precios';
\end{lstlisting}

\section{Implementación de Roles (MySQL 8.0+)}
Los roles agrupan permisos y estandarizan la seguridad, evitando la gestión individualizada que propicia errores humanos.

\subsection{Flujo de Trabajo con Roles}
\textbf{Crear y configurar el rol:}
\begin{lstlisting}
CREATE ROLE 'rol_lectura';
GRANT SELECT ON erp_empresa.* TO 'rol_lectura';
\end{lstlisting}

\textbf{Asignar y activar el rol en el usuario:}
\begin{lstlisting}
GRANT 'rol_lectura' TO 'empleado_ventas'@'%';
SET DEFAULT ROLE ALL TO 'empleado_ventas'@'%';
\end{lstlisting}

\subsection*{ Verificación de Roles}
Para comprobar qué roles están asignados y si están configurados para activarse por defecto:
\begin{lstlisting}
-- Ver las relaciones entre roles y usuarios
SELECT from_user AS 'Rol', to_user AS 'Usuario_Asignado' 
FROM mysql.role_edges 
WHERE to_user = 'empleado_ventas';

-- Ver si el rol se activa automaticamente al inicio de sesion
SELECT user, default_role_user AS 'Rol_Por_Defecto' 
FROM mysql.default_roles 
WHERE user = 'empleado_ventas';
\end{lstlisting}

\newpage

\section{Ejercicios Prácticos y Solucionario}
\textbf{Escenario:} Trabajas desde la terminal MySQL como root. Existe una base de datos llamada \texttt{sistema\_academico}.

\subsection*{Ejercicio 1: Despliegue seguro de aplicación}
\begin{itemize}
    \item Crea el usuario \texttt{webapp\_front} limitado a la IP \texttt{192.168.50.100}.
    \item Otórgale solo permisos de lectura (\texttt{SELECT}) sobre \texttt{sistema\_academico}.
    \item \textbf{Comprobación:} Escribe las consultas para verificar la creación del usuario y sus permisos a nivel de base de datos.
\end{itemize}

\textbf{Solución Ejercicio 1:}
\begin{lstlisting}
CREATE DATABASE IF NOT EXISTS sistema_academico;
CREATE USER 'webapp_front'@'192.168.50.100' IDENTIFIED BY 'PasswordFuerte123!';
GRANT SELECT ON sistema_academico.* TO 'webapp_front'@'192.168.50.100';

-- Consultas de verificacion
SELECT host, user FROM mysql.user WHERE user = 'webapp_front';
SELECT host, user, db, Select_priv, Insert_priv FROM mysql.db WHERE user = 'webapp_front';
\end{lstlisting}

\subsection*{Ejercicio 2: Diseño basado en Roles (RBAC)}
\begin{itemize}
    \item Crea los roles \texttt{auditor\_academico} (solo SELECT) y \texttt{gestor\_matriculas} (INSERT, UPDATE).
    \item Crea el usuario local \texttt{secretaria\_juan} y asígnale ambos roles, activándolos por defecto.
    \item \textbf{Comprobación:} Verifica en el diccionario de datos que los roles se han vinculado correctamente al usuario.
\end{itemize}

\textbf{Solución Ejercicio 2:}
\begin{lstlisting}
CREATE ROLE 'auditor_academico', 'gestor_matriculas';
GRANT SELECT ON sistema_academico.* TO 'auditor_academico';
GRANT INSERT, UPDATE ON sistema_academico.* TO 'gestor_matriculas';

CREATE USER 'secretaria_juan'@'localhost' IDENTIFIED BY 'PassSecre456!';
GRANT 'auditor_academico', 'gestor_matriculas' TO 'secretaria_juan'@'localhost';
SET DEFAULT ROLE ALL TO 'secretaria_juan'@'localhost';

-- Consultas de verificacion
SELECT from_user, to_user FROM mysql.role_edges WHERE to_user = 'secretaria_juan';
SELECT user, default_role_user FROM mysql.default_roles WHERE user = 'secretaria_juan';
\end{lstlisting}

\subsection*{Ejercicio 3: Auditoría y corrección de incidentes}
\begin{itemize}
    \item \textbf{Simula este error crítico:} \texttt{CREATE USER 'dev\_junior'@'\%' IDENTIFIED BY '1234'; GRANT ALL PRIVILEGES ON *.* TO 'dev\_junior'@'\%';}
    \item \textbf{Corrige el problema:} revoca los privilegios globales, asígnale solo DML (SELECT, INSERT, UPDATE, DELETE) confinado a \texttt{sistema\_academico} y fuerza la expiración de su contraseña.
    \item \textbf{Comprobación:} Demuestra mediante una consulta a \texttt{mysql.user} que el usuario ya no tiene privilegios globales.
\end{itemize}

\textbf{Solución Ejercicio 3:}
\begin{lstlisting}
-- 1. Simulacion del error
CREATE USER 'dev_junior'@'%' IDENTIFIED BY '1234';
GRANT ALL PRIVILEGES ON *.* TO 'dev_junior'@'%';

-- 2. Correccion
REVOKE ALL PRIVILEGES, GRANT OPTION FROM 'dev_junior'@'%';
GRANT SELECT, INSERT, UPDATE, DELETE ON sistema_academico.* TO 'dev_junior'@'%';
ALTER USER 'dev_junior'@'%' PASSWORD EXPIRE;

-- 3. Consultas de verificacion (Auditoria rigurosa)
-- Comprobamos que a nivel global ya no tiene poder
SELECT user, host, Select_priv, Drop_priv, Super_priv 
FROM mysql.user WHERE user = 'dev_junior';

-- Comprobamos que la contrasena esta marcada como expirada
SELECT user, host, password_expired 
FROM mysql.user WHERE user = 'dev_junior';

-- Comprobamos que sus permisos estan confinados a la BD correcta
SELECT host, user, db, Select_priv, Delete_priv, Drop_priv 
FROM mysql.db WHERE user = 'dev_junior';
\end{lstlisting}

\end{document}